% Generated draft: requires mathematical and source review.
\documentclass[11pt]{article}
\usepackage[margin=25mm]{geometry}
\usepackage{fontspec}
\setmainfont{Latin Modern Roman}
\usepackage{amsmath,amssymb,mathtools,unicode-math}
\setmathfont{Latin Modern Math}
\usepackage{enumitem,needspace,xurl,hyperref}
\setlength{\parindent}{0pt}
\setlength{\parskip}{5pt}
\setlength{\emergencystretch}{3em}
\newcommand{\N}{\mathbb{N}}
\newcommand{\Z}{\mathbb{Z}}
\newcommand{\Q}{\mathbb{Q}}
\newcommand{\R}{\mathbb{R}}
\newcommand{\C}{\mathbb{C}}
\newcommand{\F}{\mathbb{F}}
\newcommand{\A}{\mathbb{A}}
\newcommand{\PP}{\mathbb{P}}
\newcommand{\E}{\mathbb{E}}
\newcommand{\eps}{\varepsilon}
\newcommand{\dd}{\,\mathrm{d}}
\begin{document}
\section*{UnsolvedMath ID: 30003305}
\textbf{Source problem code:} catalogue entry
\subsection{Definitions and mathematical statement}
The supplied record asks about local rings without explicitly specifying commutativity, Noetherianity, or the meaning of multiplicity. The following formulation uses the customary setting of commutative Noetherian local rings and Hilbert--Samuel multiplicity; these conventions require confirmation against the cited report.

For a nonzero commutative Noetherian local ring \((A,\mathfrak m)\) of Krull dimension \(d\), define its Hilbert--Samuel multiplicity by
\[
e(A)=e_{\mathfrak m}(A)=\lim_{r\to\infty}\frac{d!\,\operatorname{length}_A(A/\mathfrak m^{r+1})}{r^d}.
\]
For a proper ideal \(J\) of a commutative Noetherian local ring \(T\), let \(\operatorname{grade}_T(J)\) be the maximal length of a \(T\)-regular sequence contained in \(J\). Here a sequence \(x_1,\ldots,x_s\) is \(T\)-regular if multiplication by \(x_i\) is injective on \(T/(x_1,\ldots,x_{i-1})\) for every \(i\), and the final quotient is nonzero. Call \(J\) perfect if
\[
\operatorname{pd}_T(T/J)=\operatorname{grade}_T(J)<\infty,
\]
where \(\operatorname{pd}_T\) denotes projective dimension.

\textbf{Question under these explicit conventions.} For every nonzero commutative Noetherian local ring \((T,\mathfrak n)\) and every proper perfect ideal \(J\subsetneq T\), must
\[
e_{\mathfrak n}(T)\leq e_{\mathfrak n/J}(T/J)
\]
hold? The two multiplicities are computed in their respective Krull dimensions. The zero ideal is included whenever it satisfies the stated definition; no positive-height restriction is supplied. No Cohen--Macaulay, equidimensionality, characteristic, or residue-field hypothesis is stated in the record. Whether the report imposes any additional hypotheses remains unresolved.
\subsection{Short English statement}
Does quotienting a local ring by a perfect ideal always weakly increase its Hilbert--Samuel multiplicity? The formulation above makes conventional ring and multiplicity assumptions explicit, subject to source confirmation. The supplied curation labels the question open; that status has not been independently verified.
\subsection{Sources}
\begin{itemize}
\item Local UnsolvedMath registry, supplied complete record: permanent ID \(30003305\), source problem code \texttt{catalogue entry}, title \emph{Multiplicity Inequality for Quotients by Perfect Ideals}, in \emph{Oberwolfach Reports Open Problems}. The record gives proposed year \(2016\) and identifies no individual proposer.
\item Supplied report citation: \emph{Asymptotic Phenomena in Local Algebra and Singularity Theory} (2017), Oberwolfach Reports, pp.\ 3199--3259, DOI \texttt{10.4171/owr/2016/57}. \url{https://doi.org/10.4171/owr/2016/57}. The background describes these reports as documenting mathematical workshops and their open-problem discussions. The supplied literature-source label is ``OWR 57 (2017).''
\item Supplied literature reference: Ma, \emph{Lech's conjecture in dimension three}. \url{https://arxiv.org/abs/1609.00095}.
\item Literature assessment dated \texttt{2026-08-21}: the supplied curation labels the problem open and describes the perfect-ideal inequality as a Lech-type strengthening known only in special cases and unresolved in general. This is supplied dated curation derived from the OWR record, not an independently verified assessment; the status and cited claims require independent checking.
\item The registry labels the statement \texttt{corrected\_verified}, but its accompanying verification note describes only safe notation, typography, and TeX-delimiter cleanup, with hypotheses, quantifiers, and subquestions preserved. This label does not establish mathematical correctness or independent source verification.
\end{itemize}
\end{document}
